\documentclass[12pt,aspectratio=169]{beamer}
\input{header_beamer}
\usetheme{Boadilla}
\usecolortheme{beaver}
\setbeamertemplate{page number in head/foot}{}
\setbeamertemplate{navigation symbols}{}
\title{Lecture-10: Correlation}
\date{}

\begin{document}
\pagestyle{empty}
\maketitle

\begin{frame}{Correlation}
Let $X:\Omega\to\R$ and $Y:\Omega\to\R$ be random variables defined on the same probability space $(\Omega,\sF,P)$. 


\begin{block}{Exercise}<2->
Show that the function $g: \R^2 \to \R$ defined by $g: (x,y) \mapsto xy$ is a Borel measurable function. 
\end{block}
\end{frame}
\begin{frame}
\begin{Definition}[Correlation] 
For two random variables $X,Y$ defined on the same probability space, 
the \textbf{correlation} between these two random variables is defined as $\E[XY]$. 
If $\E[XY] = \E[X]\E[Y]$, then the random variables $X,Y$ are called \textbf{uncorrelated}. 
\end{Definition}
\begin{block}{Lemma}<2->
If $X, Y$ are independent random variables, then they are uncorrelated. 
\end{block}
\end{frame}
\begin{frame}
\begin{Proof}
It suffices to show for $X,Y$ simple and independent random variables.
We can write $X = \sum_{x\in \sX}x\Ind{A_X(x)}$ and  $Y = \sum_{y\in \sY}y\Ind{A_Y(y)}$. 
Therefore, 
\begin{align*}
\E[XY] 
&= \sum_{(x,y) \in \sX\times\sY}xyP\set{A_X(x)\cap A_Y(y)}\\
&=  \sum_{x \in \sX}xP(A_X(x)) \sum_{y \in \sY}yP(A_Y(y))\\ 
&= \E[X]\E[Y].
\end{align*}
\end{Proof}
\end{frame}
\begin{frame}
\begin{Proof}
If $X, Y$ are independent random variables, 
then the joint distribution $F_{X,Y}(x,y) = F_X(x)F_Y(y)$ for all $(x,y) \in \R^2$.
Therefore, 
\begin{align*}
\E[XY] 
&= \int_{(x,y) \in \R^2} x y \, dF_{X,Y}(x,y) \\
&= \int_{x \in \R} x \, dF_X(x) \int_{y \in \R} y \, dF_Y(y) \\
&= \E[X] \E[Y].
\end{align*}
\end{Proof}
\end{frame}
\begin{frame}


\begin{Example}[Uncorrelated dependent random variables] 
\begin{itemize}
\item <1-> Let $X: \Omega \to \R$ be a continuous random variable with even density function $f_X:\R \to \R_+$, 
and $g: \R \to \R_+$ be another even function that is increasing for $y \in \R_+$. 
Then $g$ is Borel measurable function and $Y = g(X)$ is a random variable. 
Further, we can verify that $X, Y$ are uncorrelated and dependent random variables. 

\item <2-> To show dependence of $X$ and $Y$, 
we take positive $x,y$ such that $F_X(x) < 1$ and $x > x_y$ where $\set{x_y} = g^{-1}(y)\cap \R_+$. 
Then, we can write the set 
\EQ{
A_Y(y) = Y^{-1}(-\infty, y] 
= X^{-1}[-x_y, x_y].
} % = A_{x_y}\setminus A_{-x_y}$. 
\item <3-> Hence, we can write the joint distribution at $(x,y)$ as  
\begin{align}
F_{X,Y}(x,y) = P\set{X\le x, Y\le y} = P(A_X(x)\cap A_Y(y)) %= P(A_{x\wedge x_y}\setminus A_{-x_y}) 
= P(A_Y(y)) \notag \\ = F_Y(y) \neq F_X(x)F_Y(y).\notag
\end{align}
\end{itemize}
\end{Example}
\end{frame}
\begin{frame}
\begin{Example}[Uncorrelated dependent random variables] 
\begin{itemize}
\item <1-> Since $X$ has even density function, we have $f_X(x) = f_X(-x)$ for all $x \in \R$. 
Therefore, we have
\begin{align}
\E Xg(X)\SetIn{X < 0} 
= \int_{x < 0}xg(x)f_X(x)dx 
= \int_{u > 0}(-u)g(-u)f_X(-u)du \notag \\
= - \E Xg(X)\SetIn{X > 0}. \notag
\end{align}\\
\item <2-> The last equality follows from the fact that $g$ and $f_X$ are even. 
Therefore, we have
\begin{align}
\E[Xg(X)] 
= \E[Xg(X)\SetIn{X < 0}] + \E[Xg(X)\SetIn{X > 0}] \notag \\
= -\E[Xg(X)\SetIn{X > 0}] + \E[Xg(X)\SetIn{X > 0}]  = 0. \notag
\end{align}\\
\end{itemize}
\end{Example}
\end{frame}

\begin{frame}
\begin{block}{Theorem : AM greater than GM} 
For any two random variables $X,Y$, the correlation is upper bounded by the average of the second moments, 
with equality iff $X=Y$ almost surely. 
That is,
\EQ{
\E[XY] \le \frac{1}{2}(\E X^2 + \E Y^2). 
}
\end{block}
\end{frame}
\begin{frame}
\begin{Proof} 
This follows from the linearity and monotonicity of expectations and the fact that $(X-Y)^2 \ge 0$ with equality iff $X = Y$. 
\end{Proof}
\end{frame}
\begin{frame}
\begin{block}{Theorem : Cauchy-Schwarz inequality}  
For any two random variables $X,Y$, the correlation of absolute values of $X$ and $Y$ is upper bounded by the square root of product of second moments, 
with equality iff $X = \alpha Y$ for any constant $\alpha \in \R$. 
That is, 
\EQ{
\E\abs{XY} \le \sqrt{\E X^2\E Y^2}. 
}
\end{block}
\end{frame}
\begin{frame}
\begin{Proof} 
For two random variables $X$ and $Y$, we can define normalized random variables $W \triangleq \frac{\abs{X}}{\sqrt{\E X^2}}$ and $Z \triangleq \frac{\abs{Y}}{\sqrt{\E Y^2}}$, to get the result.
\end{Proof}
\end{frame}

\begin{frame}{Covariance}
\begin{Definition}[Covariance] 
For two random variables $X,Y$ defined on the same probability space, 
the \textbf{covariance} between these two random variables is defined as $\cov(X,Y) \triangleq \E(X-\E X)(Y-\E Y)$. 
\end{Definition}
\begin{block}{Lemma}<2->
If the random variables $X,Y$ are uncorrelated, then the covariance is zero. 
\end{block}
\end{frame}
\begin{frame}
\begin{Proof}
We can write the covariance of uncorrelated random variables $X, Y$ as 
\EQ{
\cov(X,Y) = \E(X-\E X)(Y - \E Y) = \E XY - (\E X)(\E Y) = 0.
}
\end{Proof}
\end{frame}
\begin{frame}
\begin{block}{Lemma}
Let $X: \Omega \to \R^n$ be an uncorrelated random vector and $a = (a_1, \dots, a_n) \in \R^n$, then 
\EQ{
\Var\left({\sum_{i=1}^na_iX_i}\right) = \sum_{i=1}^na_i^2\Var\left(X_i\right). 
}
\end{block}
\end{frame}
\begin{frame}
\begin{Proof}
From the linearity of expectation, we can write the variance of the linear combination as 
\EQ{
\E\left(\sum_{i=1}^na_i(X_i - \E X_i)\right)^2 = \sum_{i=1}^na_i^2\Var{X_i} + \sum_{i \neq j}a_ia_j\cov(X_i, X_j). 
}
\end{Proof}
\end{frame}
\begin{frame}
\begin{Definition}[Correlation coefficient] 
The ratio of covariance of two random variables $X,Y$ and the square root of product of their variances is called the \textbf{correlation coefficient} and denoted by
\EQ{
\rho_{X,Y} \triangleq \frac{\cov(X,Y)}{\sqrt{\Var(X), \Var(Y)}}. 
}
\end{Definition}
\begin{block}{Theorem : Correlation coefficient}<2->
For any two random variables $X,Y$, the absolute value of correlation coefficient is less than or equal to unity, with equality iff $X = \alpha Y + \beta$ almost surely for constants $\alpha = \sqrt{\frac{\Var(X)}{\Var(Y)}}$ and $\beta = \E X - \alpha \E Y$. 
\end{block}
\end{frame}
\begin{frame}
\begin{Proof} 
\begin{itemize}
\item <1-> For two random variables $X$ and $Y$, we can define normalized random variables $W \triangleq \frac{X-\E X}{\sqrt{\Var(X)}}$ and $Z \triangleq \frac{Y-\E Y}{\sqrt{\Var(Y)}}$. 
Applying the AM-GM inequality to random variables $W, Z$, we get 
\EQ{
\abs{\cov(X,Y)}  \le \sqrt{\Var(X) \Var(Y)}. 
}
\item <2-> Recall that equality is achieved iff $W = Z$ almost surely or equivalently iff $X = \alpha Y + \beta$ almost surely.  
Taking $U = -Y$, we see that $-\cov(X,Y) \le \sqrt{\Var(X)\Var(Y)}$, and hence the result follows.
\end{itemize}
\end{Proof}
\end{frame}


\begin{frame}{$L^p$ spaces}



\begin{Definition}
A pair $(p,q) \in \R^2$ where $p, q \ge 1$ and $\frac{1}{p}+\frac{1}{q} = 1$, i
s called the \textbf{conjugate pair}, 
and the spaces $L^p$ and $L^q$ are called \textbf{dual spaces}. 
\end{Definition}


\begin{Example}<2->
The dual of $L^1$ space is $L^\infty$. 
The space $L^2$ is dual of itself, and called a \textbf{Hilbert space}. 
\end{Example}
\end{frame}

\begin{frame}
\begin{block}{Theorem : H\"older's inequality} 
Consider a conjugate pair $(p,q)$ and random variables $X \in L^p, Y \in L^q$. 
Then, 
\EQ{
\E\abs{XY} \le \norm{X}_p\norm{Y}_q.%(\E\abs{X}^p)^{\frac{1}{p}}(\E\abs{Y}^q)^{\frac{1}{q}}.
}
\end{block}
\end{frame}
\begin{frame}
\begin{Proof}
\begin{itemize}
\item <1-> Consider a random variable $Z: \Omega \to \set{p\ln v, q \ln w}$ with probability mass function $\set{\frac{1}{p}, \frac{1}{q}}$. 
It follows from Jensen's inequality applied to the convex function $f(x) = e^x$ and the random variable $Z$, that 
\EQ{
vw 
= f(\E Z) 
\le \E f(Z)  
= \frac{v^p}{p} + \frac{w^q}{q}. 
}
\item <2-> It follows that for any random variables $V,W$, we have $VW \le \frac{V^p}{p} + \frac{W^q}{q}$. 
Taking expectation on both sides, 
we get from the monotonicity of expectation that  
$
\E VW \le \frac{\E V^p}{p} +\frac{\E W^q}{q}.
$
Taking $V \triangleq \frac{\abs{X}}{\norm{X}_p}$ and $W \triangleq \frac{\abs{Y}}{\norm{Y}_q}$, we get the result.
\end{itemize}
\end{Proof}
\end{frame}
\begin{frame}
\begin{Definition}
For a pair of random variables $(X,Y) \in (L^p,L^q)$ for conjugate pair $(p,q)$, 
we can define inner product $\inner{}:L^p\times L^q \to \R$ by 
\EQ{
\inner{}(X,Y) \triangleq \inner{X,Y} \triangleq \E XY.
}
\end{Definition}
\begin{block}{Reminder}<2->
For $X \in L^p$ and $Y\in L^q$, the expectation $\E\abs{XY}$ is finite from H\"older's inequality.
Therefore, the inner product $\inner{X,Y} = \E[XY]$ is well defined and finite.
\end{block}
\begin{block}{Reminder}<3->
This inner product is well defined for the conjugate pair $(1,\infty)$.
\end{block}
\end{frame}

\begin{frame}
\begin{block}{Theorem : Minkowski's inequality} 
For $1 \le p < \infty$, let $X,Y\in L^p$ be two random variables defined on a probability space $(\Omega, \sF, P)$. 
Then, 
\EQ{
\norm{X+ Y}_p \le \norm{X}_p + \norm{Y}_p,
} 
with inequality iff $X = \alpha Y$ for some $\alpha \ge 0$ or $Y = 0$. 
\end{block}
\end{frame}

\begin{frame}
\begin{Proof}
\begin{itemize}
\item <1-> Since addition is a Borel measurable function, $X+Y$ is a random variable. 
We first show that $X+Y \in L^p$, when $X,Y \in L^p$.  
To this end, we observe that $g:\R_+\to\R_+$ defined by $g(x)= x^p$ for all $x \in \R_+$, 
is a convex function for $p \ge 1$. 
From the convexity of $g$, we have 
\begin{align}
\abs{\frac{1}{2}X + \frac{1}{2}Y}^p 
\le \abs{\frac{1}{2}\abs{X} + \frac{1}{2}\abs{Y}}^p
= g(\frac{1}{2}\abs{X} + \frac{1}{2}\abs{Y})\notag \\
\le \frac{1}{2}g(\abs{X}) + \frac{1}{2}g(\abs{Y}) = \frac{1}{2}\abs{X}^p + \frac{1}{2}\abs{Y}^p. \notag
\end{align}\\
\item <2-> This implies that $\abs{X+Y}^p \le 2^{p-1}(\abs{X}^p + \abs{Y}^p)$. 
The inequality holds trivially if $\norm{X+Y}_p = 0$. 
Therefore, we assume that $\norm{X+Y}_p > 0$, without any loss of generality. 
\end{itemize}
\end{Proof}
\end{frame}

\begin{frame}
\begin{Proof}[Proof continued.]
\begin{itemize}
\item <1-> Using the definition of $\norm{}_p$, triangle inequality, 
and linearity of expectation 
we get 
\begin{align}
\norm{X+Y}_p^p 
= \E[\abs{X+Y}\abs{X+Y}^{p-1}] \le \E([\abs{X}+\abs{Y})\abs{X+Y}^{p-1}] \notag \\
= \E\abs{X}\abs{X+Y}^{p-1} +  \E\abs{Y}\abs{X+Y}^{p-1}. \notag
\end{align}\\
\item <2-> From the H\"older's inequality applied to conjugate pair $(p,q)$ to the two products on RHS, we get 
\EQ{
\norm{X+Y}_p^p \le (\norm{X}_p+\norm{Y}_p)\norm{\abs{X+Y}^{p-1}}_q
}
\item <3-> Recall that $q = \frac{p}{p-1}$. %, and hence $(p-1)q = p$. 
Therefore, $\norm{\abs{X+Y}^{p-1}}_q = (\E\abs{X+Y}^p)^{1-\frac{1}{p}}$ and the result follows.
\end{itemize}
\end{Proof}
\end{frame}
\begin{frame}
\begin{block}{Reminder}
We have shown that the map $\norm{}_p$ is a norm by proving the Minkowski's inequality. 
Therefore, $L^p$ is a normed vector space.    
We can define distance between two random variables $X_1, X_2 \in L^p$ by the norm $\norm{X_1-X_2}_p$. 
\end{block}
\end{frame}
\end{document}
